\section{The ramified first norm and the three cubic unit classes}
\label{rsx-section}

Retain \(\zeta^2-\zeta+1=0\),
\(\mathcal O=\mathbb Z[\zeta]\), and
\(N(x+y\zeta)=x^2+xy+y^2\). For positive coprime integers
\(a,b\), put
\[
 c=a+b,\quad d=a-b,\quad M=a^2+ab+b^2,\quad
 F=a^4+3a^3b+5a^2b^2+3ab^3+b^4.
\]
This section treats \(M=3U^2\), retaining its ramified coefficient.
Its proof does not apply the preceding theorem with \(M=U^2\).

\begin{proposition}[The exact ramified common-source models]
\label{rsx-models}
Put
\begin{align}
 A_0&=S^3-3ST^2-T^3,& B_0&=3ST(S+T),\nonumber\\
 C_0&=A_0+B_0=S^3+3S^2T-T^3.
 \label{rsx-cubics}
\end{align}
The coordinates \((A,B)\) of
\(\epsilon(S+T\zeta)^3\), for
\(\epsilon=1,\zeta,\zeta^2\), respectively, are
\begin{equation}
 (A_0,B_0),\qquad(-B_0,C_0),\qquad(-C_0,A_0).
 \label{rsx-unit-coordinates}
\end{equation}
Let \(L_{\epsilon,3}:[S:T]\mapsto[A:B]\), and let
\(\mathcal K\) be the smooth projective conic with the indicated
morphism to \(\mathbb P^1_r\):
\begin{equation}
 3V^2+W^2=12L^2,\qquad
 r=\frac{V^2-3L^2}{3L^2}.
 \label{rsx-conic}
\end{equation}
The normalization \(D^{\rm ram}_{3,\epsilon}\) of their
projective fiber product is geometrically connected of genus six.
Its three genus-two quotients are the smooth projective models of
\begin{align}
 (H^{\rm ram}_1):\quad Y_1^2&=3B(A+B),\nonumber\\
 (H^{\rm ram}_2):\quad Y_2^2&=3B(B-3A),\nonumber\\
 (H^{\rm ram}_3):\quad Y_3^2&=(A+B)(B-3A).
 \label{rsx-quotients}
\end{align}
These equations all use the same projective source \([S:T]\).
\end{proposition}
\begin{proof}
Multiplication with \(\zeta^2=\zeta-1\) proves
\eqref{rsx-cubics}--\eqref{rsx-unit-coordinates}.
Over \(\mathbb Q(\zeta)\), the fractional linear map
\[
 \rho(x)=\frac{x+\zeta}{x+\bar\zeta}
\]
gives the identity of projective rational maps
\begin{equation}
 \rho(L_{\epsilon,3}(s))=\epsilon^2\rho(s)^3,\qquad s=S/T.
 \label{rsx-conjugacy}
\end{equation}
Thus \(L_{\epsilon,3}\) is a degree-three morphism, its
coordinates have no common projective zero, and its branch
values are \(-\zeta,-\bar\zeta\). Each of
\(\infty,-1,1/3\) therefore has three simple preimages,
with the three fibers pairwise disjoint.

The conic in \eqref{rsx-conic} is smooth. When \(L=0\),
its coordinate \(V\) cannot vanish, so the two homogeneous
coordinates defining \(r\) never vanish simultaneously.
On \(L=1\), its equations become
\begin{equation}
 v^2=3(1+r),\qquad w^2=3(1-3r).
 \label{rsx-affine-squares}
\end{equation}
The square classes of \(3(A+B)/B\) and \(3(B-3A)/B\)
are independent over the geometric function field: each has
three simple zeros absent from the other. The fiber product
therefore has geometric degree four over \(\mathbb P^1_{[S:T]}\).
Its nine branch points are the three fibers above
\(\infty,-1,1/3\), with inertia two at each. At their
common poles the ratio of the two functions is a geometric
local unit, so the inertia is still two. Riemann--Hurwitz gives
\[
 2g(D^{\rm ram}_{3,\epsilon})-2=4(-2)+9\cdot2=10.
\]
On \(B\ne0\), set \(Y_1=Bv\), \(Y_2=Bw\), and
\(Y_3=Bvw/3\). These give \eqref{rsx-quotients}.
Each right side has six distinct projective roots, proving
the stated genera. In particular the third quotient is
unchanged by the twist, while the first two are twisted by three.

The normalization has morphisms to both the source projective
line and the whole conic. Their morphisms to \(\mathbb P^1_r\)
agree on a dense open chart and hence everywhere. Thus no
affine presentation discards exceptional or infinite fibers.
\end{proof}

\begin{theorem}[Complete local exclusion of all three ramified covers]
\label{rsx-local-empty}
For each \(\epsilon\in\{1,\zeta,\zeta^2\}\),
\[
 D^{\rm ram}_{3,\epsilon}(\mathbb Q_3)=\varnothing .
\]
\end{theorem}
\begin{proof}
A point of \(\mathcal K(\mathbb Q_3)\) with \(L=0\)
would give \(W^2=-3V^2\) with \(V\ne0\), impossible by
valuation parity. Every such point is therefore affine.
In
\[
 3v^2+w^2=12
\]
the two nonzero terms on the left have valuations of opposite
parity. They cannot cancel at the minimum, which must equal
\(v_3(12)=1\). It must be the valuation of \(3v^2\).
Also \(v=0\) would require a square of odd valuation.
Allowing \(w=0\), this proves
\begin{equation}
 v\in\mathbb Z_3^\times,\qquad w\in3\mathbb Z_3,
 \qquad v_3(r)=v_3((v^2-3)/3)=-1.
 \label{rsx-required-depth}
\end{equation}

Now normalize any \([S:T]\in\mathbb P^1(\mathbb Q_3)\)
to \(S,T\in\mathbb Z_3\), at least one a unit. This includes
\(T=0\). If \(S\not\equiv T\pmod3\), then
\[
 A_0\equiv C_0\equiv(S-T)^3\not\equiv0\pmod3.
\]
One of \(S,T,S+T\) vanishes modulo three: if both coordinates
are nonzero their distinct residues are opposite. Consequently
\(B_0\in9\mathbb Z_3\). The three ratios in
\eqref{rsx-unit-coordinates} have the following exhaustive
possibilities:
\begin{equation}
\begin{array}{c|c}
\epsilon& r=A/B\\ \hline
1&v_3(r)\le-2,\ \text{or }r=\infty\\
\zeta&v_3(r)\ge2,\ \text{or }r=0\\
\zeta^2&v_3(r)=0 .
\end{array}
\label{rsx-unit-depths}
\end{equation}

If \(S\equiv T\pmod3\), both are units. Write \(S=T+3K\).
The exact identities
\begin{align}
 A_0&=-3T^3+27K^2T+27K^3,\nonumber\\
 B_0&=3ST(S+T),\nonumber\\
 C_0&=3(T^3+9KT^2+18K^2T+9K^3)
 \label{rsx-diagonal-expansion}
\end{align}
show that all three have valuation one: after division by three
their residues are \(-T^3,2T^3,T^3\), respectively.
Each of the three ratios is then a unit.

No projective source for any of the three units has finite ratio
of valuation \(-1\). A rational point on the normalization would
have everywhere defined images on that source line and on
\(\mathcal K\), with the same \(r\). This contradicts
\eqref{rsx-required-depth}. The proof therefore includes every
branch fiber, infinite source, and ramified input.
\end{proof}

\begin{theorem}[An actual integer cube forces an unramified first norm]
\label{rsx-integer-gate}
Let \(a,b,Q\) be positive integers with \(\gcd(a,b)=1\).
Then
\begin{equation}
 F=Q^3\quad\Longrightarrow\quad3\nmid M.
 \label{rsx-main-implication}
\end{equation}
In the exact integral cubic extraction
\begin{equation}
 ab+M\zeta=\epsilon(S+T\zeta)^3,\qquad
 \epsilon\in\{1,\zeta,\zeta^2\},
 \label{rsx-extraction}
\end{equation}
the identity unit \(\epsilon=1\) cannot occur.
\end{theorem}
\begin{proof}
The actual element \(z=ab+M\zeta\) satisfies
\begin{equation}
 N(z)=(ab)^2+abM+M^2=F,\quad
 \gcd(ab,M)=1,\quad F\equiv(a^2+b^2)^2\equiv1\pmod3.
 \label{rsx-primitive-norm}
\end{equation}
Indeed a prime dividing \(ab,M\) would divide both \(a,b\).
For the last congruence their squared sum modulo three is
one or two, because they are not both zero.

The ring \(\mathcal O\) is norm Euclidean: rounding the two
real coordinates in the basis \(1,\zeta\) gives an error
of norm at most \(3/4<1\). It is therefore a UFD; its units
are the six powers of \(\zeta\). An inert rational prime in
\(N(z)\) would divide \(z\) as a rational integer, contrary
to primitivity. At a split prime, both conjugate prime
factors cannot occur, since their product would divide both
coordinates. The ramified prime above three is absent by
\eqref{rsx-primitive-norm}. Every occurring orientation
thus has exponent equal to its rational exponent in \(N(z)\).

Under \(F=Q^3\) all those exponents are multiples of three.
This proves \eqref{rsx-extraction}, with integral \(S,T\),
\(\gcd(S,T)=1\), and \(N(S+T\zeta)=Q\). A sign is absorbed
into the cube root to leave the three displayed units.
Since \(3\nmid Q\),
\begin{equation}
 S\not\equiv T\pmod3 .
 \label{rsx-unramified-root}
\end{equation}
This is an equality of actual integral elements. It neither
divides a projective output by its content nor drops a
nonunit residual.

A primitive pair always has \(v_3(M)\in\{0,1\}\).
Use \(M=(a-b)^2+3ab\). If \(a\not\equiv b\pmod3\),
then \(M\) is a unit. Otherwise \(a,b\) are units, the
second term has valuation one, and the first has valuation
at least two.

By \eqref{rsx-unramified-root}, \(A_0,C_0\) are units
and \(B_0\in9\mathbb Z_3\). Thus the second coordinate in
\eqref{rsx-unit-coordinates} is either zero or of valuation
at least two for \(\epsilon=1\), and is a unit for each of
the other two units. It cannot have valuation one.
Since this coordinate is the positive integer \(M\)
of valuation at most one, the identity unit is impossible
and the other two force \(3\nmid M\).
\end{proof}

\begin{corollary}[The actual ramified square/cube family is empty]
\label{rsx-square-cube}
There are no positive integers \(a,b,U,Q\) with \(\gcd(a,b)=1\)
and
\[
 M=3U^2,\qquad F=Q^3.
\]
\end{corollary}
\begin{proof}
Theorem~\ref{rsx-integer-gate} immediately contradicts \(3\mid M\).
The full local curve obstruction gives a second presentation
of the same implication, with its source checks as follows.
The exact extraction \eqref{rsx-extraction} gives the source
\([S:T]\) with \(r=ab/M\). Set \(v=c/U,w=d/U\).
Since \(c^2=M+ab\) and \(d^2=M-3ab\), these satisfy
\eqref{rsx-affine-squares} at that same source.
Positivity gives
\[
 0<r\le1/3,\quad v>0,\quad v>|w|,
 \quad\text{equivalently}\quad \sqrt3<v\le2 .
\]
Equality \(r=1/3\) would force \(a=b=1,U=1,F=13\), which
is not a cube. Hence \(0<r<1/3\), avoiding all three branch
targets; this actual point lifts to the smooth normalization
and contradicts Theorem~\ref{rsx-local-empty}.
No converse for arbitrary rational points is used.
\end{proof}

\paragraph{Other elementary local checks.}
The polynomial \(F\) is odd at each of the three nonzero pairs
modulo two. Modulo five, its values \(F(S,1)\) for
\(S=0,1,2,3,4\) are \(1,3,2,2,1\); when \(5\mid b\),
one has \(F\equiv a^4\not\equiv0\).
Thus \(2,3,5\nmid Q\) in an actual cubic profile.
The ramified square profile would also have \(U\) odd and
a three-adic unit. These are supplementary checks, not a
claim of a complete list of local conditions.

\begin{corollary}[Numerical and projective propagation]
\label{rsx-propagation}
For positive integers \(a,b,Q,g\), with \(\gcd(a,b)=1\) and
\(3\mid g\), the equality \(F=Q^g\) implies \(3\nmid M\).
In particular \(M=3R^h,F=Q^g\) is impossible for positive
integers \(R,h\) under these assumptions.
For every \(g=3k\), \(k\) a positive integer, and
\(\epsilon\in\{1,\zeta,\zeta^2\}\), the analogous normalized
ramified cover \(D^{\rm ram}_{g,\epsilon}\) has no
\(\mathbb Q_3\)-point.
\end{corollary}
\begin{proof}
Numerically rewrite \(Q^g=(Q^{g/3})^3\) and apply
Theorem~\ref{rsx-integer-gate}; there is no parity restriction.
For the geometric assertion let \(L_{\epsilon,g}\) be the
coordinate ratio of \(\epsilon(S+T\zeta)^g\). Then
\[
 L_{\epsilon,g}=L_{\epsilon,3}\circ L_{1,k},
 \qquad \rho(L_{\epsilon,g}(s))=\epsilon^2\rho(s)^g .
\]
Each of \(\infty,-1,1/3\) has \(g\) simple preimages,
with the three fibers disjoint. The two square classes
remain independent by their disjoint simple zeros, so the
fiber-product normalization is geometrically connected,
including when \(g\) is even. Retaining the conic coordinates
and applying \(L_{1,k}\) to the source gives a nonconstant
rational map \(D^{\rm ram}_{g,\epsilon}\to
D^{\rm ram}_{3,\epsilon}\), retaining \(r\).
It extends to the smooth projective models by
\cite[Lemma~53.2.2 and Theorem~53.2.6]{StacksCurves2026}.
A \(\mathbb Q_3\)-point would contradict
Theorem~\ref{rsx-local-empty}.
This geometric statement concerns these specified
coefficients; it does not identify arbitrary unit classes
modulo \(g\)-th powers with the three cubic classes.
The numerical assertion needs no such identification.
\end{proof}

\paragraph{Why the global extraction cannot be replaced by a numerical
local test.}
The two numerical norm equations alone can be soluble over
\(\mathbb Q_3\). At \(a=1,b=4\) one has \(M=21,F=541\).
The equation \(U^2=7\) has a simple root modulo three and
hence a root in \(\mathbb Q_3\). To solve \(Q^3=541\), put
\(Q=1+3t\); the resulting equation is
\[
 t+3t^2+3t^3=60.
\]
Its reduction has the simple root \(t=0\) modulo three,
so Hensel's lemma supplies \(Q\in\mathbb Q_3\).
Thus local numerical solubility does not contradict
Theorem~\ref{rsx-local-empty}. The latter concerns exact
common-source covers with one of three global units.
Their relevance to integer solutions comes from the
primitive global UFD extraction, not from the numerical
local norm alone.

\paragraph{Ordinary proof and remaining scope.}
The complete argument is ordinary mathematics; no Lean
formalization of the UFD extraction, curve construction,
or local exclusion is claimed. It is a fixed exact-profile
result. It gives no uniform height estimate for varying
residual coefficients and no theorem putting every possible
ABC counterexample in these perfect-power families.
Nonunit second residuals can change the cubic-unit equations;
their coverage, other indices, and uniform point-height
bounds remain open. The preceding unramified theorem
retains its separate proof and hypotheses.
